\section{总结与延伸}

本讲的主线可以压缩成一句话：AI 能力增长不是单一模型曲线，而是一条由 context、compute、capital 与 culture 共同推动的全栈转型。Context 决定系统接触什么真实任务和反馈，compute 决定可承受的实验规模，capital 决定基础设施以多快速度建设，culture 与 institutions 决定目标、责任和资源分配。任何一层被忽略，另一层的扩张都可能变成浪费或风险。

\begin{figure}[H]
\centering
\includegraphics[width=0.92\textwidth]{images/01-03-45-food-for-thought.jpg}
\caption{课程留给学生的两个问题：如何实现和平的 compute transition，以及个人要怎样参与。}
\figsource{Stanford Online 官方视频 01:03:45，`images/01-03-45-food-for-thought.jpg`。}
\end{figure}

\begin{importantbox}{读图：结尾不是旁观式预测题}
第一问要求把“和平转型”拆成可及性、互操作、市场集中、能源、跨境供应链、安全与公共利益；第二问要求学生提出自己能验证的行动，例如定义资源标准、做可移植 workload、公开 benchmark 方法、研究调度与结算、参与标准组织或解释政策后果。课程项目的价值不只在展示一个模型，而在让学生进入真实基础设施对话。
\end{importantbox}

\begin{knowledgebox}{课堂提示：把自己当成 active participant}
讲者最后强调，学生不必等到成为 CEO、研究主管或监管者才有资格参与。清楚写出接口痛点、公开可复现实验、记录迁移失败、提出标准草案、向行业讲者追问证据，都可能改变讨论。Preparedness 的终点不是更会预测谁会赢，而是能够在不确定时期提出更好的问题并建设可检查的系统。
\end{knowledgebox}

\subsection{八条可迁移结论}

\begin{enumerate}
  \item 先画全栈依赖，再讨论单层突破；模型、产品和算力都嵌在更大的资源系统中。
  \item 把 scaling law 当经验模型，明确观测区间、外推假设和失效条件。
  \item 把 context 拆成任务、工具、环境、权限与反馈，不要只理解为长 prompt。
  \item 审计 RL 时同时检查 environment、reward、verifier 和责任边界。
  \item 阅读产业案例时跟踪入口、模型、执行环境、数据回流与分发控制权。
  \item 对宏观图表区分技术、采用、资本和价格信号，不用一类证据替代另一类结论。
  \item 使用历史类比时迁移机制，不迁移确定时间表；特别检查资产寿命和标准化条件。
  \item Compute 要成为广泛可访问基础设施，需要技术标准与执行制度共同成熟。
\end{enumerate}

\subsection{拓展阅读}

继续学习时，建议从三条可复现路径展开。第一，访问 Stanford CS153 课程官网与 Stanford Online 官方视频，对照后续嘉宾怎样修改本讲的四约束框架。第二，选择一个 coding agent 或训练任务，画出其 context flow、执行环境、验证信号与数据授权。第三，比较两个云或 GPU 市场的 job specification、SLA 和计费口径，实际检查 workload 能否迁移，而不是只比较标价。

\begin{knowledgebox}{讲义提醒：后续课程怎样使用本讲}
每听一位嘉宾，至少记录四件事：他所在的系统层级、他认为最稀缺的资源、他展示的证据类型、他希望由市场还是制度解决的问题。十周结束后，把这些答案放在同一张表里，冲突之处往往比共识更有学习价值。
\end{knowledgebox}

\subsection{实践作业：审计一条真实 Compute 路径}

为了把本讲的宏观框架变成系统能力，可以选择自己实际使用的一项 AI 服务，从一次请求出发追踪到计算资源。作业重点不是猜测供应商的全部内部架构，而是明确哪些事实可观察、哪些来自公开文档、哪些仍是推断。最终交付应让另一位读者能够复查资源接口、context 回流与责任边界。

\begin{enumerate}
  \item 画出请求从客户端、应用服务、模型 API 到推理集群的路径，标出身份、数据和工具权限在哪一层生效。
  \item 选择一次可重复任务，记录延迟、失败类型、输出质量与重试行为；说明 reward 或成功条件能否被自动验证。
  \item 比较两个供应渠道的规格、价格和迁移步骤，检查 ``GPU-hour''、token price 或 API 名称是否掩盖不可替代性。
  \item 提出一个最小标准改进，例如可移植 job spec、统一故障码、性能披露或数据删除证明，并说明由什么机构推动采用。
\end{enumerate}

\begin{importantbox}{验收标准}
一份合格作业应包含来源链接、可复现实验步骤、至少一个失败案例、事实与推断的明确分栏，以及对安全和公共利益的影响说明。若结论只能写成“某公司更强”或“需要更多算力”，说明分析还没有进入 Frontier Systems 所要求的接口、机制与制度层。
\end{importantbox}

\end{document}
